\chapter{Polar Equations of Conic Sections (Pages 1--14)}
\label{ch:polar_conics}

% =========================================================================
% PAGE 01
% =========================================================================
\section{Page 1: Focus-Directrix Definition and Standard Polar Equations}
\label{sec:page01}

\begin{conceptbox}[Manuscript Overview: Page 1]
Page 1 introduces the polar coordinate system applied to conic sections. The focus $S$ of the conic is taken as the pole (origin), and the principal axis of symmetry passing through $S$ perpendicular to the directrix is selected as the initial line.
\end{conceptbox}

\subsection{Geometric Definition and Polar Setup}
A conic section is the locus of a point $P$ moving in a plane such that the ratio of its distance from a fixed point $S$ (the \textbf{focus}) to its perpendicular distance from a fixed straight line $ZM$ (the \textbf{directrix}) is a constant $e$ (the \textbf{eccentricity}):
\begin{equation}
    \frac{SP}{PM} = e \iff SP = e \cdot PM.
\end{equation}
The nature of the conic is governed by the value of $e$:
\begin{itemize}[noitemsep]
    \item $e = 1$: \textbf{Parabola}
    \item $0 \le e < 1$: \textbf{Ellipse} ($e=0$ represents a circle with directrix at infinity)
    \item $e > 1$: \textbf{Hyperbola}
\end{itemize}

\begin{center}
\begin{tikzpicture}[scale=1.15, >=Stealth]
    % Initial line
    \draw[->, thick] (-3,0) -- (6.5,0) node[right] {Initial Line $SX$};
    \draw[dashed, gray!60] (0,-3.2) -- (0,3.2);

    % Pole S
    \coordinate (S) at (0,0);
    \node[below left] at (S) {$S \text{ (Focus/Pole)}$};
    \fill (S) circle (2pt);

    % Directrix ZM
    \draw[very thick, primary] (-2.4,-3) -- (-2.4,3) node[above] {Directrix $ZM$};
    \coordinate (Z) at (-2.4,0);
    \node[below left] at (Z) {$Z$};
    \fill (Z) circle (1.5pt);

    % Semi-latus rectum LSL'
    \coordinate (L) at (0,2.1);
    \coordinate (Lprime) at (0,-2.1);
    \draw[thick, secondary] (Lprime) -- (L) node[above] {$L(l, \pi/2)$};
    \node[below right] at (Lprime) {$L'(l, -\pi/2)$};
    \fill (L) circle (1.5pt);
    \fill (Lprime) circle (1.5pt);
    \draw[dashed, gray] (L) -- (-2.4,2.1) node[midway, above] {$d$};

    % Arbitrary Point P(r, theta)
    \coordinate (P) at (36:3.4);
    \draw[very thick, accent] (S) -- (P) node[midway, above left] {$r$};
    \node[above right] at (P) {$P(r, \theta)$};
    \fill (P) circle (2pt);

    % Projections
    \coordinate (N) at ($(S)!(P)!(6,0)$);
    \draw[dashed] (P) -- (N);
    \node[below] at (N) {$N$};

    \coordinate (M) at (-2.4, {3.4*sin(36)});
    \draw[dashed, secondary] (P) -- (M) node[midway, above] {$PM$};
    \node[left] at (M) {$M$};
    \fill (M) circle (1.5pt);

    % Angle theta
    \draw[->, thick, primary] (0.8,0) arc (0:36:0.8) node[midway, right=2pt] {$\theta$};

    % Distance SZ = d
    \draw[<->] (-2.4,-0.5) -- (0,-0.5) node[midway, below] {$SZ = d = l/e$};
\end{tikzpicture}
\end{center}

\subsection{Step-by-Step Derivation of the Standard Equation}
Let $Z$ be the foot of the perpendicular from the focus $S$ to the directrix, and let $SZ = d$.
By definition of the semi-latus rectum $l = SL$:
The point $L$ on the conic has vectorial angle $\theta = \frac{\pi}{2}$. Its distance from the directrix is equal to $SZ = d$.
Applying the focus-directrix definition to $L$:
\begin{equation}
    SL = e \cdot ZM_L \implies l = e \cdot d \iff d = \frac{l}{e}.
\end{equation}
Now, let $P(r, \theta)$ be any point on the conic. From $P$, draw $PN \perp SX$ and $PM \perp ZM$.
From the geometry of the figure:
\begin{align}
    SN &= r\cos\theta, \\
    PM &= ZN = SZ + SN = d + r\cos\theta.
\end{align}
Using $SP = e \cdot PM$:
\begin{align}
    r &= e(d + r\cos\theta) \nonumber\\
    r &= ed + er\cos\theta \nonumber\\
    r &= l + er\cos\theta \nonumber\\
    r(1 - e\cos\theta) &= l \quad \text{(when directrix is to the right of the focus)}.
\end{align}
When the directrix is chosen to the left of the focus (as standard in Indian university curricula):
\begin{equation}
    \boxed{\frac{l}{r} = 1 + e\cos\theta}
\end{equation}

\subsection{Standard Orientations and Rotations}
\begin{theorembox}[Orientations of the Polar Conic]
\begin{enumerate}
    \item \textbf{Directrix to the left of the focus}:
    \begin{equation}
        \frac{l}{r} = 1 + e\cos\theta
    \end{equation}
    \item \textbf{Directrix to the right of the focus}:
    \begin{equation}
        \frac{l}{r} = 1 - e\cos\theta
    \end{equation}
    \item \textbf{Major axis inclined at an angle $\alpha$ to the initial line}:
    \begin{equation}
        \frac{l}{r} = 1 + e\cos(\theta - \alpha)
    \end{equation}
    \item \textbf{Directrix parallel to the initial line}:
    \begin{equation}
        \frac{l}{r} = 1 \pm e\sin\theta
    \end{equation}
\end{enumerate}
\end{theorembox}

% =========================================================================
% PAGE 02
% =========================================================================
\newpage
\section{Page 2: Focal Chords and the Semi-Latus Rectum}
\label{sec:page02}

\begin{conceptbox}[Manuscript Overview: Page 2]
Page 2 introduces the concept of a focal chord (a chord passing through the focus $S$) and establishes one of the central classical theorems of analytical geometry: the semi-latus rectum is the harmonic mean of the segments of any focal chord.
\end{conceptbox}

\subsection{Geometric Formulation}
Let $PSQ$ be a focal chord of the conic $\frac{l}{r} = 1 + e\cos\theta$, passing through the pole $S$.
Let the vectorial angle of $P$ be $\alpha$. Since $P, S, Q$ are collinear in that order, the radius vector $SQ$ is directed in the opposite sense, meaning the vectorial angle of $Q$ is $\alpha + \pi$.

\begin{center}
\begin{tikzpicture}[scale=1.1, >=Stealth]
    \coordinate (S) at (0,0);
    \node[below left] at (S) {$S \text{ (Focus)}$};
    \fill (S) circle (2pt);

    \draw[->, thick] (-1,0) -- (5.5,0) node[right] {Initial Line};

    % Focal Chord PSQ
    \coordinate (P) at (38:3.6);
    \coordinate (Q) at (218:1.8);
    \draw[very thick, primary] (Q) -- (P);
    \node[above right] at (P) {$P(SP, \alpha)$};
    \node[below left] at (Q) {$Q(SQ, \alpha+\pi)$};
    \fill (P) circle (2pt);
    \fill (Q) circle (2pt);

    % Segment labels
    \draw[dashed, secondary] (S) -- (P) node[midway, above left] {$SP$};
    \draw[dashed, secondary] (S) -- (Q) node[midway, below right] {$SQ$};

    % Angles
    \draw[->, thick, ForestGreen] (0.8,0) arc (0:38:0.8) node[midway, right] {$\alpha$};
    \draw[->, thick, accent] (0.5,0) arc (0:218:0.5) node[midway, left] {$\alpha+\pi$};
\end{tikzpicture}
\end{center}

\begin{theorembox}[Theorem: Harmonic Mean Property of Focal Chords]
The semi-latus rectum of any conic is the harmonic mean between the segments of any focal chord:
\begin{equation}
    \boxed{\frac{1}{SP} + \frac{1}{SQ} = \frac{2}{l} = \text{constant}}
\end{equation}
\end{theorembox}

\begin{proof}
Since $P(SP, \alpha)$ lies on the conic $\frac{l}{r} = 1 + e\cos\theta$:
\begin{equation}
    \frac{l}{SP} = 1 + e\cos\alpha. \label{eq:p2_P}
\end{equation}
Similarly, since $Q(SQ, \alpha + \pi)$ lies on the conic:
\begin{align}
    \frac{l}{SQ} &= 1 + e\cos(\alpha + \pi) \nonumber\\
                 &= 1 - e\cos\alpha. \label{eq:p2_Q}
\end{align}
Adding equations \eqref{eq:p2_P} and \eqref{eq:p2_Q} side by side:
\begin{align}
    \frac{l}{SP} + \frac{l}{SQ} &= (1 + e\cos\alpha) + (1 - e\cos\alpha) \nonumber\\
                                &= 2.
\end{align}
Dividing both sides by $l$:
\begin{equation}
    \frac{1}{SP} + \frac{1}{SQ} = \frac{2}{l}.
\end{equation}
Since $l$ is a fixed constant for a given conic, the sum of the reciprocals of the focal segments $SP$ and $SQ$ is constant and completely independent of the inclination $\alpha$ of the focal chord.
\end{proof}

\begin{notebox}[Pedagogical Significance]
This result proves that regardless of how a focal chord is rotated about the focus, the harmonic mean of its two parts $SP$ and $SQ$ remains strictly invariant and equal to the semi-latus rectum $l$.
\end{notebox}

% =========================================================================
% PAGE 03
% =========================================================================
\newpage
\section{Page 3: Equation of a Chord Joining Points $(\alpha \pm \beta)$}
\label{sec:page03}

\begin{conceptbox}[Manuscript Overview: Page 3]
Page 3 derives the general equation of a secant straight line joining two points on the conic $\frac{l}{r} = 1 + e\cos\theta$ whose vectorial angles are chosen symmetrically as $\alpha - \beta$ and $\alpha + \beta$.
\end{conceptbox}

\subsection{Algebraic Derivation}
The general equation of a straight line in polar coordinates can be written in the linear harmonic form:
\begin{equation}
    \frac{l}{r} = A\cos\theta + B\sin\theta. \label{eq:p3_genline}
\end{equation}
Equivalently, it can be re-parameterized as:
\begin{equation}
    \frac{l}{r} = e\cos\theta + C\cos(\theta - \delta).
\end{equation}
Let this straight line pass through the two points on the conic:
\begin{equation}
    P\left(r_1, \; \alpha - \beta\right) \quad \text{and} \quad Q\left(r_2, \; \alpha + \beta\right).
\end{equation}
Since $P$ lies on the conic:
\begin{equation}
    \frac{l}{r_1} = 1 + e\cos(\alpha - \beta).
\end{equation}
Since $P$ also lies on the chord \eqref{eq:p3_genline}:
\begin{equation}
    \frac{l}{r_1} = A\cos(\alpha - \beta) + B\sin(\alpha - \beta).
\end{equation}
Equating the two expressions for $\frac{l}{r_1}$:
\begin{equation}
    A\cos(\alpha - \beta) + B\sin(\alpha - \beta) = 1 + e\cos(\alpha - \beta) \iff (A - e)\cos(\alpha - \beta) + B\sin(\alpha - \beta) = 1. \label{eq:p3_eq1}
\end{equation}
Similarly, for point $Q(r_2, \alpha + \beta)$:
\begin{equation}
    (A - e)\cos(\alpha + \beta) + B\sin(\alpha + \beta) = 1. \label{eq:p3_eq2}
\end{equation}

Subtracting equation \eqref{eq:p3_eq2} from equation \eqref{eq:p3_eq1}:
\begin{align}
    (A - e)[\cos(\alpha - \beta) - \cos(\alpha + \beta)] + B[\sin(\alpha - \beta) - \sin(\alpha + \beta)] &= 0 \nonumber\\
    (A - e)[2\sin\alpha\sin\beta] - B[2\cos\alpha\sin\beta] &= 0.
\end{align}
Assuming $\beta \neq 0$ (so that $\sin\beta \neq 0$), dividing by $2\sin\beta$ yields:
\begin{equation}
    (A - e)\sin\alpha - B\cos\alpha = 0 \implies \frac{A - e}{\cos\alpha} = \frac{B}{\sin\alpha} = k.
\end{equation}
Thus:
\begin{equation}
    A - e = k\cos\alpha \implies A = e + k\cos\alpha, \quad \text{and} \quad B = k\sin\alpha.
\end{equation}
Substitute these back into equation \eqref{eq:p3_eq1}:
\begin{align}
    (k\cos\alpha)\cos(\alpha - \beta) + (k\sin\alpha)\sin(\alpha - \beta) &= 1 \nonumber\\
    k[\cos\alpha\cos(\alpha - \beta) + \sin\alpha\sin(\alpha - \beta)] &= 1 \nonumber\\
    k\cos[\alpha - (\alpha - \beta)] &= 1 \nonumber\\
    k\cos\beta &= 1 \implies k = \sec\beta.
\end{align}
Therefore:
\begin{equation}
    A = e + \sec\beta\cos\alpha, \quad B = \sec\beta\sin\alpha.
\end{equation}
Substituting $A$ and $B$ into the line equation $\frac{l}{r} = A\cos\theta + B\sin\theta$:
\begin{align}
    \frac{l}{r} &= (e + \sec\beta\cos\alpha)\cos\theta + (\sec\beta\sin\alpha)\sin\theta \nonumber\\
    &= e\cos\theta + \sec\beta(\cos\theta\cos\alpha + \sin\theta\sin\alpha) \nonumber\\
    &= e\cos\theta + \sec\beta\cos(\theta - \alpha).
\end{align}

\begin{theorembox}[Equation of a Chord in Symmetric Polar Form]
The straight line joining the points on $\frac{l}{r} = 1 + e\cos\theta$ with vectorial angles $\alpha - \beta$ and $\alpha + \beta$ is:
\begin{equation}
    \boxed{\frac{l}{r} = e\cos\theta + \sec\beta\cos(\theta - \alpha)}
\end{equation}
\end{theorembox}

% =========================================================================
% PAGE 04
% =========================================================================
\newpage
\section{Page 4: Alternate Chord Form and the Tangent at Point $\alpha$}
\label{sec:page04}

\begin{conceptbox}[Manuscript Overview: Page 4]
Page 4 rewrites the chord equation for two arbitrary vectorial angles $\alpha$ and $\beta$, derives the equation of the tangent line as the limiting case, and establishes the condition for tangency of a general line.
\end{conceptbox}

\subsection{Chord Joining Vectorial Angles $\alpha$ and $\beta$}
Let the two endpoints be given directly as $P(\alpha)$ and $Q(\beta)$.
Their mean angle is $\frac{\alpha + \beta}{2}$ and their half-difference is $\frac{\alpha - \beta}{2}$.
Setting the mean angle as $\alpha'$ and half-difference as $\beta'$ in the formula from Page 3:
\begin{equation}
    \boxed{\frac{l}{r} = e\cos\theta + \sec\left(\frac{\alpha - \beta}{2}\right)\cos\left(\theta - \frac{\alpha + \beta}{2}\right)}
\end{equation}

\subsection{Derivation of the Tangent at Point $\alpha$}
The tangent line to the conic at point $P(\alpha)$ is the limiting position of the secant chord joining $\alpha - \beta$ and $\alpha + \beta$ as the second point approaches the first, which corresponds to the limit $\beta \to 0$.
Since $\lim_{\beta \to 0} \sec\beta = \sec 0 = 1$:

\begin{theorembox}[Equation of the Tangent Line]
The equation of the tangent to the conic $\frac{l}{r} = 1 + e\cos\theta$ at the point with vectorial angle $\alpha$ is:
\begin{equation}
    \boxed{\frac{l}{r} = e\cos\theta + \cos(\theta - \alpha)}
\end{equation}
Expanding the cosine term:
\begin{equation}
    \boxed{\frac{l}{r} = (e + \cos\alpha)\cos\theta + \sin\alpha\sin\theta}
\end{equation}
\end{theorembox}

\subsection{Condition of Tangency for a Straight Line}
Let a straight line be given by:
\begin{equation}
    \frac{l}{r} = A\cos\theta + B\sin\theta. \label{eq:p4_testline}
\end{equation}
Comparing equation \eqref{eq:p4_testline} with the tangent equation $\frac{l}{r} = (e + \cos\alpha)\cos\theta + \sin\alpha\sin\theta$:
\begin{equation}
    A = e + \cos\alpha \implies A - e = \cos\alpha, \quad \text{and} \quad B = \sin\alpha.
\end{equation}
Squaring and adding:
\begin{equation}
    (A - e)^2 + B^2 = \cos^2\alpha + \sin^2\alpha = 1.
\end{equation}

\begin{theorembox}[Condition of Tangency]
The straight line $\frac{l}{r} = A\cos\theta + B\sin\theta$ touches the conic $\frac{l}{r} = 1 + e\cos\theta$ if and only if:
\begin{equation}
    \boxed{(A - e)^2 + B^2 = 1}
\end{equation}
\end{theorembox}

% =========================================================================
% PAGE 05
% =========================================================================
\newpage
\section{Page 5: Condition of Tangency and the Normal Line}
\label{sec:page05}

\begin{conceptbox}[Manuscript Overview: Page 5]
Page 5 expresses the tangency condition in perpendicular distance form $r\cos(\theta-\phi) = p$ and derives the polar equation of the normal line to the conic at point $\alpha$.
\end{conceptbox}

\subsection{Tangency Condition in Standard Form $r\cos(\theta - \phi) = p$}
Let the straight line be given in perpendicular normal form:
\begin{equation}
    r\cos(\theta - \phi) = p \iff \frac{p}{r} = \cos\theta\cos\phi + \sin\theta\sin\phi.
\end{equation}
Multiplying by $\frac{l}{p}$:
\begin{equation}
    \frac{l}{r} = \left(\frac{l}{p}\cos\phi\right)\cos\theta + \left(\frac{l}{p}\sin\phi\right)\sin\theta.
\end{equation}
Setting $A = \frac{l}{p}\cos\phi$ and $B = \frac{l}{p}\sin\phi$ into the condition $(A - e)^2 + B^2 = 1$:
\begin{align}
    \left(\frac{l}{p}\cos\phi - e\right)^2 + \left(\frac{l}{p}\sin\phi\right)^2 &= 1 \nonumber\\
    \frac{l^2}{p^2}\cos^2\phi - \frac{2le}{p}\cos\phi + e^2 + \frac{l^2}{p^2}\sin^2\phi &= 1 \nonumber\\
    \frac{l^2}{p^2}(\cos^2\phi + \sin^2\phi) - \frac{2le}{p}\cos\phi + e^2 &= 1 \nonumber\\
    \frac{l^2}{p^2} - \frac{2le}{p}\cos\phi + e^2 &= 1.
\end{align}
Multiplying through by $p^2$:
\begin{equation}
    \boxed{(l\cos\phi - ep)^2 + l^2\sin^2\phi = p^2 \iff l^2 - 2elp\cos\phi + p^2(e^2 - 1) = 0}
\end{equation}

\subsection{Derivation of the Equation of the Normal Line}
The normal line at point $P(\alpha)$ is the straight line through $P$ perpendicular to the tangent at $P$.
The equation of the tangent at $P(\alpha)$ is:
\begin{equation}
    \frac{l}{r} = (e + \cos\alpha)\cos\theta + \sin\alpha\sin\theta.
\end{equation}
In Cartesian coordinates, setting $x = r\cos\theta, y = r\sin\theta$:
\begin{equation}
    (e + \cos\alpha)x + (\sin\alpha)y = l.
\end{equation}
The slope of this tangent line is:
\begin{equation}
    m_T = -\frac{e + \cos\alpha}{\sin\alpha}.
\end{equation}
The slope of the normal line must satisfy $m_N \cdot m_T = -1$:
\begin{equation}
    m_N = \frac{\sin\alpha}{e + \cos\alpha}.
\end{equation}
The equation of any line perpendicular to the tangent is:
\begin{equation}
    -(\sin\alpha)x + (e + \cos\alpha)y = K.
\end{equation}
Converting to polar coordinates:
\begin{equation}
    r[-\sin\alpha\cos\theta + (e + \cos\alpha)\sin\theta] = K \implies r[e\sin\theta + \sin(\theta - \alpha)] = K.
\end{equation}
Since this line passes through $P(r_\alpha, \alpha)$ on the conic, where $r_\alpha = \frac{l}{1 + e\cos\alpha}$:
\begin{equation}
    K = r_\alpha [e\sin\alpha + \sin(\alpha - \alpha)] = \frac{l \cdot e\sin\alpha}{1 + e\cos\alpha}.
\end{equation}

\begin{theorembox}[Equation of the Normal Line in Polar Form]
\begin{equation}
    \boxed{\frac{le\sin\alpha}{1 + e\cos\alpha} \cdot \frac{1}{r} = e\sin\theta + \sin(\theta - \alpha)}
\end{equation}
\end{theorembox}

% =========================================================================
% PAGE 06
% =========================================================================
\newpage
\section{Page 6: Chord of Contact and the Polar Line}
\label{sec:page06}

\begin{conceptbox}[Manuscript Overview: Page 6]
Page 6 derives the equation of the chord of contact of tangents drawn from an external point $T(r_1, \theta_1)$ to the conic, and defines the polar of a point with respect to a conic.
\end{conceptbox}

\subsection{Derivation of the Chord of Contact}
Let $T(r_1, \theta_1)$ be an external point from which two tangents are drawn to the conic $\frac{l}{r} = 1 + e\cos\theta$, touching the curve at points $P(\alpha)$ and $Q(\beta)$.
The equation of the tangent at any point with vectorial angle $\phi$ is:
\begin{equation}
    \frac{l}{r} = e\cos\theta + \cos(\theta - \phi).
\end{equation}
Since both tangents pass through $T(r_1, \theta_1)$:
\begin{align}
    \frac{l}{r_1} - e\cos\theta_1 &= \cos(\theta_1 - \alpha), \label{eq:p6_tang1}\\
    \frac{l}{r_1} - e\cos\theta_1 &= \cos(\theta_1 - \beta). \label{eq:p6_tang2}
\end{align}
Now consider the equation of the chord joining $\alpha$ and $\beta$:
\begin{equation}
    \frac{l}{r} = e\cos\theta + \sec\left(\frac{\alpha - \beta}{2}\right)\cos\left(\theta - \frac{\alpha + \beta}{2}\right).
\end{equation}
From equations \eqref{eq:p6_tang1} and \eqref{eq:p6_tang2}, since $\cos(\theta_1 - \alpha) = \cos(\theta_1 - \beta)$, we have:
\begin{equation}
    \theta_1 - \alpha = -(\theta_1 - \beta) \implies \theta_1 = \frac{\alpha + \beta}{2}.
\end{equation}
Moreover:
\begin{equation}
    \frac{l}{r_1} - e\cos\theta_1 = \cos\left(\frac{\alpha - \beta}{2}\right) \implies \sec\left(\frac{\alpha - \beta}{2}\right) = \frac{1}{\frac{l}{r_1} - e\cos\theta_1}.
\end{equation}
Substituting these into the chord equation:
\begin{align}
    \frac{l}{r} - e\cos\theta &= \frac{\cos(\theta - \theta_1)}{\frac{l}{r_1} - e\cos\theta_1}.
\end{align}
Cross-multiplying yields:

\begin{theorembox}[Equation of the Chord of Contact and Polar Line]
The equation of the chord of contact of tangents from $T(r_1, \theta_1)$ (and the \textbf{polar} of $T$) with respect to $\frac{l}{r} = 1 + e\cos\theta$ is:
\begin{equation}
    \boxed{\left(\frac{l}{r} - e\cos\theta\right)\left(\frac{l}{r_1} - e\cos\theta_1\right) = \cos(\theta - \theta_1)}
\end{equation}
\end{theorembox}

% =========================================================================
% PAGE 07
% =========================================================================
\newpage
\section{Page 7: Classification of Conics and Analysis of Eccentricity}
\label{sec:page07}

\begin{conceptbox}[Manuscript Overview: Page 7]
Page 7 deals with transforming arbitrary polar equations into the standard canonical form $\frac{l'}{r} = 1 \pm e'\cos\theta$ to determine the eccentricity, semi-latus rectum, and geometric nature of the conic.
\end{conceptbox}

\subsection{Canonical Reduction Method}
Given an equation of the form:
\begin{equation}
    \frac{A}{r} = B + C\cos\theta + D\sin\theta,
\end{equation}
we reduce it to the standard canonical form:
\begin{enumerate}
    \item Divide throughout by the constant term $B$ ($B \neq 0$):
    \begin{equation}
        \frac{A/B}{r} = 1 + \frac{C}{B}\cos\theta + \frac{D}{B}\sin\theta.
    \end{equation}
    \item Combine the cosine and sine terms using $R = \sqrt{(C/B)^2 + (D/B)^2}$ and $\tan\phi = D/C$:
    \begin{equation}
        \frac{A/B}{r} = 1 + R\cos(\theta - \phi).
    \end{equation}
    \item Identify the parameters:
    \begin{equation}
        \text{Semi-latus rectum } l' = \frac{A}{B}, \quad \text{Eccentricity } e' = R = \frac{\sqrt{C^2 + D^2}}{|B|}.
    \end{equation}
\end{enumerate}

\subsection{Detailed Analysis of the Manuscript Example}
The manuscript analyzes the equation:
\begin{equation}
    \frac{l}{r} = 4 - 5\cos\theta.
\end{equation}
Dividing throughout by $4$:
\begin{equation}
    \frac{l/4}{r} = 1 - \frac{5}{4}\cos\theta.
\end{equation}
Comparing with $\frac{l'}{r} = 1 - e\cos\theta$:
\begin{itemize}[noitemsep]
    \item Semi-latus rectum: $l' = \frac{l}{4}$
    \item Eccentricity: $e = \frac{5}{4} = 1.25$
\end{itemize}

\begin{notebox}[Correction of Manuscript Slip]
In the manuscript notes on Page 7, the student wrote:
\begin{equation*}
    \text{``Since } e = \frac{5}{4} \dots \text{ ellipse''} \quad \text{\textbf{[ERROR]}}
\end{equation*}
By the fundamental definition of conic sections:
\begin{itemize}[noitemsep]
    \item If $e < 1$, the conic is an \textbf{ellipse}.
    \item If $e = 1$, the conic is a \textbf{parabola}.
    \item If $e > 1$, the conic is a \textbf{hyperbola}.
\end{itemize}
Since $e = \frac{5}{4} = 1.25 > 1$, this conic is unambiguously a \textbf{hyperbola}.
\end{notebox}

% =========================================================================
% PAGE 08
% =========================================================================
\newpage
\section{Page 8: Asymptotes of a Hyperbola in Polar Coordinates}
\label{sec:page08}

\begin{conceptbox}[Manuscript Overview: Page 8]
Page 8 derives the equations of the asymptotes of a hyperbola in polar coordinates. An asymptote is defined analytically as a tangent line whose point of contact lies at infinity.
\end{conceptbox}

\subsection{Mathematical Derivation}
Let the hyperbola have equation:
\begin{equation}
    \frac{l}{r} = 1 + e\cos\theta, \quad (e > 1).
\end{equation}
The point of contact $(\infty, \alpha)$ lies at infinity as $r \to \infty$. In this limit, $\frac{l}{r} \to 0$:
\begin{equation}
    1 + e\cos\alpha = 0 \implies \cos\alpha = -\frac{1}{e}.
\end{equation}
Since $\sin^2\alpha = 1 - \cos^2\alpha$:
\begin{equation}
    \sin\alpha = \pm \sqrt{1 - \left(-\frac{1}{e}\right)^2} = \pm \sqrt{1 - \frac{1}{e^2}} = \pm \frac{\sqrt{e^2 - 1}}{e}.
\end{equation}
The general equation of the tangent at point $\alpha$ is:
\begin{align}
    \frac{l}{r} &= e\cos\theta + \cos(\theta - \alpha) \nonumber\\
                &= e\cos\theta + \cos\theta\cos\alpha + \sin\theta\sin\alpha \nonumber\\
                &= (e + \cos\alpha)\cos\theta + (\sin\alpha)\sin\theta.
\end{align}
Substituting $\cos\alpha = -\frac{1}{e}$ and $\sin\alpha = \pm \sqrt{1 - \frac{1}{e^2}}$:
\begin{align}
    \frac{l}{r} &= \left(e - \frac{1}{e}\right)\cos\theta \pm \sqrt{1 - \frac{1}{e^2}}\sin\theta.
\end{align}
Transposing the cosine term:
\begin{equation}
    \frac{l}{r} - \left(e - \frac{1}{e}\right)\cos\theta = \pm \sqrt{1 - \frac{1}{e^2}}\sin\theta.
\end{equation}
Squaring both sides to combine both asymptotes into a single second-degree equation:

\begin{theorembox}[Combined Equation of Asymptotes]
The combined polar equation of the pair of asymptotes of the hyperbola $\frac{l}{r} = 1 + e\cos\theta$ is:
\begin{equation}
    \boxed{\left[\frac{l}{r} - \left(e - \frac{1}{e}\right)\cos\theta\right]^2 = \left(1 - \frac{1}{e^2}\right)\sin^2\theta}
\end{equation}
\end{theorembox}

% =========================================================================
% PAGE 09
% =========================================================================
\newpage
\section{Page 9: Problem 1 --- Tangent Segment Intercepted by Directrix}
\label{sec:page09}

\begin{examplebox}[Problem 1 (Manuscript Page 9)]
\textbf{Problem Statement:}
Prove that the portion of any tangent to a conic intercepted between the curve and the directrix subtends a right angle at the corresponding focus.
\tcblower
\textbf{Proof:}
Let the focus $S$ be taken as the pole, and let the conic be:
\begin{equation}
    \frac{l}{r} = 1 + e\cos\theta.
\end{equation}
The directrix corresponding to the focus $S$ is the line $r\cos\theta = -\frac{l}{e} \iff \frac{l}{r} = -e\cos\theta$.
Let $P(\alpha)$ be any point of contact on the conic. The equation of the tangent at $P$ is:
\begin{equation}
    \frac{l}{r} = e\cos\theta + \cos(\theta - \alpha). \label{eq:p9_tan}
\end{equation}
Let this tangent intersect the directrix at point $K(r_K, \theta_K)$.
At the directrix, $\frac{l}{r_K} = -e\cos\theta_K$.
Substituting this into equation \eqref{eq:p9_tan}:
\begin{align}
    -e\cos\theta_K &= e\cos\theta_K + \cos(\theta_K - \alpha) \nonumber\\
    \cos(\theta_K - \alpha) &= -2e\cos\theta_K.
\end{align}
If we use the standard directrix $\frac{l}{r} = e\cos\theta$:
\begin{align}
    e\cos\theta_K &= e\cos\theta_K + \cos(\theta_K - \alpha) \nonumber\\
    \cos(\theta_K - \alpha) &= 0.
\end{align}
Hence:
\begin{equation}
    \theta_K - \alpha = \pm \frac{\pi}{2} \iff |\angle KSP| = \frac{\pi}{2} = 90^\circ.
\end{equation}
Thus, the intercepted segment $PK$ subtends a right angle ($90^\circ$) at the focus $S$.
\end{examplebox}

% =========================================================================
% PAGE 10
% =========================================================================
\newpage
\section{Page 10: Problem 2 --- Sum of Reciprocal Products of Focal Chords}
\label{sec:page10}

\begin{examplebox}[Problem 2 (Manuscript Page 10)]
\textbf{Problem Statement:}
If $PSP'$ and $QSQ'$ are two mutually perpendicular focal chords of a conic $\frac{l}{r} = 1 + e\cos\theta$, prove that:
\begin{equation*}
    \frac{1}{SP \cdot SP'} + \frac{1}{SQ \cdot SQ'} = \text{constant}.
\end{equation*}
\tcblower
\textbf{Solution:}
Let the vectorial angle of $P$ be $\alpha$. Since $PSP'$ is a straight line through the pole $S$, the vectorial angle of $P'$ is $\alpha + \pi$.
From the equation of the conic:
\begin{align}
    \frac{l}{SP} &= 1 + e\cos\alpha, \\
    \frac{l}{SP'} &= 1 + e\cos(\alpha + \pi) = 1 - e\cos\alpha.
\end{align}
Multiplying these two relations:
\begin{align}
    \frac{l^2}{SP \cdot SP'} &= (1 + e\cos\alpha)(1 - e\cos\alpha) = 1 - e^2\cos^2\alpha \nonumber\\
    \frac{1}{SP \cdot SP'} &= \frac{1 - e^2\cos^2\alpha}{l^2}. \label{eq:p10_1}
\end{align}
Since $QSQ'$ is perpendicular to $PSP'$, the vectorial angle of $Q$ is $\alpha + \frac{\pi}{2}$, and that of $Q'$ is $\alpha + \frac{3\pi}{2}$.
Hence:
\begin{align}
    \frac{l}{SQ} &= 1 + e\cos\left(\alpha + \frac{\pi}{2}\right) = 1 - e\sin\alpha, \\
    \frac{l}{SQ'} &= 1 + e\cos\left(\alpha + \frac{3\pi}{2}\right) = 1 + e\sin\alpha.
\end{align}
Multiplying:
\begin{align}
    \frac{l^2}{SQ \cdot SQ'} &= (1 - e\sin\alpha)(1 + e\sin\alpha) = 1 - e^2\sin^2\alpha \nonumber\\
    \frac{1}{SQ \cdot SQ'} &= \frac{1 - e^2\sin^2\alpha}{l^2}. \label{eq:p10_2}
\end{align}
Adding equations \eqref{eq:p10_1} and \eqref{eq:p10_2}:
\begin{align}
    \frac{1}{SP \cdot SP'} + \frac{1}{SQ \cdot SQ'} &= \frac{(1 - e^2\cos^2\alpha) + (1 - e^2\sin^2\alpha)}{l^2} \nonumber\\
    &= \frac{2 - e^2(\cos^2\alpha + \sin^2\alpha)}{l^2} \nonumber\\
    &= \boxed{\frac{2 - e^2}{l^2} = \text{constant}}
\end{align}
This quantity is completely independent of $\alpha$, completing the proof.
\end{examplebox}

% =========================================================================
% PAGE 11
% =========================================================================
\newpage
\section{Page 11: Problem 3 --- Tangents at Extremities of a Focal Chord}
\label{sec:page11}

\begin{examplebox}[Problem 3 (Manuscript Page 11)]
\textbf{Problem Statement:}
Prove that the tangents at the ends of any focal chord of a conic intersect on the corresponding directrix.
\tcblower
\textbf{Proof:}
Let $PSQ$ be any focal chord of the conic $\frac{l}{r} = 1 + e\cos\theta$.
Let the vectorial angle of $P$ be $\alpha$. Then the vectorial angle of $Q$ is $\alpha + \pi$.
The equation of the tangent at $P(\alpha)$ is:
\begin{equation}
    \frac{l}{r} = e\cos\theta + \cos(\theta - \alpha). \label{eq:p11_tan1}
\end{equation}
The equation of the tangent at $Q(\alpha + \pi)$ is:
\begin{align}
    \frac{l}{r} &= e\cos\theta + \cos[\theta - (\alpha + \pi)] \nonumber\\
                &= e\cos\theta + \cos[(\theta - \alpha) - \pi] \nonumber\\
                &= e\cos\theta - \cos(\theta - \alpha). \label{eq:p11_tan2}
\end{align}
To find their point of intersection, subtract equation \eqref{eq:p11_tan2} from equation \eqref{eq:p11_tan1}:
\begin{equation}
    2\cos(\theta - \alpha) = 0 \implies \cos(\theta - \alpha) = 0.
\end{equation}
Substituting $\cos(\theta - \alpha) = 0$ back into equation \eqref{eq:p11_tan1}:
\begin{equation}
    \frac{l}{r} = e\cos\theta.
\end{equation}
This is precisely the equation of the directrix of the conic.
Hence, the point of intersection of the tangents at the extremities of any focal chord always lies on the directrix.
\end{examplebox}

% =========================================================================
% PAGE 12
% =========================================================================
\newpage
\section{Page 12: Problem 4 --- Perpendicular Focal Tangents of a Parabola}
\label{sec:page12}

\begin{examplebox}[Problem 4 (Manuscript Page 12)]
\textbf{Problem Statement:}
For a parabola ($e = 1$), prove that the tangents at the ends of any focal chord intersect at right angles on the directrix.
\tcblower
\textbf{Proof:}
For a parabola, $e = 1$, and its equation is:
\begin{equation}
    \frac{l}{r} = 1 + \cos\theta.
\end{equation}
Let $PSQ$ be a focal chord with $P$ at $\alpha$ and $Q$ at $\alpha + \pi$.
From Page 11, the tangents intersect on the directrix where $\cos(\theta - \alpha) = 0 \implies \theta = \alpha \pm \frac{\pi}{2}$.
Now consider the inclination of the tangents.
In Cartesian coordinates, the equation of the tangent at point $\alpha$ is:
\begin{equation}
    (1 + \cos\alpha)x + (\sin\alpha)y = l \implies 2\cos^2\left(\frac{\alpha}{2}\right)x + 2\sin\left(\frac{\alpha}{2}\right)\cos\left(\frac{\alpha}{2}\right)y = l.
\end{equation}
The slope of this tangent is:
\begin{equation}
    m_1 = -\frac{1 + \cos\alpha}{\sin\alpha} = -\frac{2\cos^2(\alpha/2)}{2\sin(\alpha/2)\cos(\alpha/2)} = -\cot\left(\frac{\alpha}{2}\right).
\end{equation}
Similarly, the slope of the tangent at $Q(\alpha + \pi)$ is obtained by replacing $\alpha$ with $\alpha + \pi$:
\begin{align}
    m_2 &= -\cot\left(\frac{\alpha + \pi}{2}\right) = -\cot\left(\frac{\alpha}{2} + \frac{\pi}{2}\right) = - \left(-\tan\frac{\alpha}{2}\right) = \tan\left(\frac{\alpha}{2}\right).
\end{align}
Evaluating the product of the slopes:
\begin{equation}
    m_1 \cdot m_2 = \left[-\cot\left(\frac{\alpha}{2}\right)\right] \cdot \left[\tan\left(\frac{\alpha}{2}\right)\right] = -1.
\end{equation}
Since $m_1 m_2 = -1$, the two tangents are mutually perpendicular ($90^\circ$).
Thus, for a parabola, the tangents at the ends of any focal chord intersect at right angles on the directrix.
\end{examplebox}

% =========================================================================
% PAGE 13
% =========================================================================
\newpage
\section{Page 13: Problem 5 --- Director Circle in Polar Coordinates}
\label{sec:page13}

\begin{examplebox}[Problem 5 (Manuscript Page 13)]
\textbf{Problem Statement:}
Find the locus of the point of intersection of two perpendicular tangents to the conic $\frac{l}{r} = 1 + e\cos\theta$ (Equation of the Director Circle).
\tcblower
\textbf{Solution:}
Let the tangents be drawn at points with vectorial angles $\alpha$ and $\beta$.
Their equations are:
\begin{align}
    \frac{l}{r} &= (e + \cos\alpha)\cos\theta + \sin\alpha\sin\theta, \\
    \frac{l}{r} &= (e + \cos\beta)\cos\theta + \sin\beta\sin\theta.
\end{align}
Their Cartesian slopes are:
\begin{equation}
    m_1 = -\frac{e + \cos\alpha}{\sin\alpha}, \quad m_2 = -\frac{e + \cos\beta}{\sin\beta}.
\end{equation}
Since the tangents are perpendicular, $m_1 m_2 = -1$:
\begin{align}
    \frac{(e + \cos\alpha)(e + \cos\beta)}{\sin\alpha\sin\beta} &= -1 \nonumber\\
    e^2 + e(\cos\alpha + \cos\beta) + \cos\alpha\cos\beta + \sin\alpha\sin\beta &= 0 \nonumber\\
    e^2 + 2e\cos\left(\frac{\alpha+\beta}{2}\right)\cos\left(\frac{\alpha-\beta}{2}\right) + \cos(\alpha - \beta) &= 0. \label{eq:p13_perp}
\end{align}
Let $(r_0, \theta_0)$ be the coordinates of their point of intersection.
Solving the two tangent equations simultaneously gives:
\begin{equation}
    \theta_0 = \frac{\alpha + \beta}{2}, \quad \text{and} \quad \frac{l}{r_0} = e\cos\theta_0 + \cos\left(\frac{\alpha - \beta}{2}\right).
\end{equation}
Let $\gamma = \frac{\alpha - \beta}{2}$. Then $\cos\gamma = \frac{l}{r_0} - e\cos\theta_0$, and $\cos(\alpha - \beta) = 2\cos^2\gamma - 1$.
Substituting into equation \eqref{eq:p13_perp}:
\begin{align}
    e^2 + 2e\cos\theta_0\left(\frac{l}{r_0} - e\cos\theta_0\right) + 2\left(\frac{l}{r_0} - e\cos\theta_0\right)^2 - 1 &= 0 \nonumber\\
    e^2 - 1 + \frac{2el}{r_0}\cos\theta_0 - 2e^2\cos^2\theta_0 + 2\left(\frac{l^2}{r_0^2} - \frac{2el}{r_0}\cos\theta_0 + e^2\cos^2\theta_0\right) &= 0 \nonumber\\
    (e^2 - 1) - \frac{2el}{r_0}\cos\theta_0 + \frac{2l^2}{r_0^2} &= 0.
\end{align}
Multiplying by $r_0^2$ and generalizing $(r_0, \theta_0) \to (r, \theta)$:
\begin{equation}
    \boxed{(e^2 - 1)r^2 - 2elr\cos\theta + 2l^2 = 0}
\end{equation}
This is the polar equation of the \textbf{director circle} of the conic.
\end{examplebox}

% =========================================================================
% PAGE 14
% =========================================================================
\newpage
\section{Page 14: Problem 6 --- Locus of the Pole of a Subtended Chord}
\label{sec:page14}

\begin{examplebox}[Problem 6 (Manuscript Page 14)]
\textbf{Problem Statement:}
Find the locus of the pole of a chord of the conic $\frac{l}{r} = 1 + e\cos\theta$ which subtends a constant angle $2\beta$ at the focus.
\tcblower
\textbf{Solution:}
Let $T(r_1, \theta_1)$ be the pole of the chord.
The chord is the chord of contact of tangents drawn from $T(r_1, \theta_1)$ to the conic.
Let the points of contact of these tangents be $P$ and $Q$, with vectorial angles $\alpha - \beta$ and $\alpha + \beta$, so that the angle subtended by the chord $PQ$ at the focus $S$ is:
\begin{equation}
    \angle PSQ = (\alpha + \beta) - (\alpha - \beta) = 2\beta = \text{constant}.
\end{equation}
The equation of the chord of contact from $T(r_1, \theta_1)$ is:
\begin{equation}
    \frac{l}{r} = e\cos\theta + \sec\beta\cos(\theta - \alpha).
\end{equation}
As derived on Page 6, the coordinates of the pole $T(r_1, \theta_1)$ satisfy:
\begin{equation}
    \theta_1 = \alpha, \quad \text{and} \quad \frac{l}{r_1} - e\cos\theta_1 = \cos\beta.
\end{equation}
Therefore:
\begin{equation}
    \frac{l}{r_1} = e\cos\theta_1 + \cos\beta.
\end{equation}
Dividing by $\cos\beta$:
\begin{equation}
    \frac{l\sec\beta}{r_1} = 1 + e\sec\beta\cos\theta_1.
\end{equation}
Generalizing to arbitrary polar coordinates $(r, \theta)$:
\begin{equation}
    \boxed{\frac{l\sec\beta}{r} = 1 + (e\sec\beta)\cos\theta}
\end{equation}
\textbf{Conclusion:}
The locus of the pole is another conic of the same type, having the same focus and axis, with new semi-latus rectum $l' = l\sec\beta$ and new eccentricity $e' = e\sec\beta$.
\end{examplebox}
